% Generated draft: requires mathematical and source review.
\documentclass[11pt]{article}
\usepackage[margin=25mm]{geometry}
\usepackage{fontspec}
\setmainfont{Latin Modern Roman}
\usepackage{amsmath,amssymb,mathtools,unicode-math}
\setmathfont{Latin Modern Math}
\usepackage{enumitem,needspace,xurl,hyperref}
\setlength{\parindent}{0pt}
\setlength{\parskip}{5pt}
\setlength{\emergencystretch}{3em}
\newcommand{\N}{\mathbb{N}}
\newcommand{\Z}{\mathbb{Z}}
\newcommand{\Q}{\mathbb{Q}}
\newcommand{\R}{\mathbb{R}}
\newcommand{\C}{\mathbb{C}}
\newcommand{\F}{\mathbb{F}}
\newcommand{\A}{\mathbb{A}}
\newcommand{\PP}{\mathbb{P}}
\newcommand{\E}{\mathbb{E}}
\newcommand{\eps}{\varepsilon}
\newcommand{\dd}{\,\mathrm{d}}
\begin{document}
\section*{UnsolvedMath ID: 20000250}
\textbf{Source problem code:} catalogue entry
\subsection{Definitions and mathematical statement}
Let $d,r$ be positive integers and $X$ a CW complex of dimension $d$, not assumed to have finitely many cells. For $1\leq i\leq r$, prescribe $c_i\in H^{2i}(X;\Z)$. For each positive integer $m$, let
\[
 \rho_m:H^{2i}(X;\Z)\longrightarrow H^{2i}(X;\Z/m\Z)
\]
be coefficient reduction. Nori's question asks whether
\[
 \rho_m(c_i)=0\quad\text{for every }m\geq1\text{ and every }1\leq i\leq r
\]
implies existence of a topological complex vector bundle $E\to X$ of rank exactly $r$ with $c_i(E)=c_i$ for all $i$. A finite-dimensional CW complex must not be replaced by a finite CW complex.

The report observes, through the coefficient long exact sequence, that the hypothesis is equivalently $c_i\in mH^{2i}(X;\Z)$ for every positive integer $m$. This is individual divisibility by every integer; no compatible inverse system of chosen roots is stipulated. In classifying-space terms the target is a lift of the specified Chern data through $BU(r)\to\prod_{i=1}^rK(\Z,2i)$.

The supplied research aid attributes an affirmative answer to a consequence of Benoist's 2026 universal-multiple lemma. It claims that for fixed $d,r$ a single multiplier $M(d,r)$ works independently of the space and cohomology tuple, and that choosing $M(d,r)$-th roots of the prescribed classes produces the exact rank-$r$ bundle. The lemma's precise scaling statement and target-skeleton proof are not included. It explicitly claims cellular approximation handles infinitely many cells without an inverse-limit or phantom-map gap; that quantifier is retained as a reported proof requirement.

Its candidate example takes, for every prime $p$, the mapping telescope
\[
 T=\operatorname{Tel}(S^5\xrightarrow{p}S^5\xrightarrow{p}\cdots).
\]
For every nonzero $\gamma\in H^6(T;\Z)$ with the report's interpretation in $\Z_p/\Z$, it claims a rank-three complex bundle with $c_1=c_2=0$ and $c_3=\gamma$, where $\Z_p$ denotes the $p$-adic integers. The exact cohomology identification and witness construction are not supplied. The affirmative theorem and example remain unreviewed source reports, not new results proved in this draft.
\subsection{Short English statement}
Can integral Chern classes vanishing modulo every positive integer be realized by a bundle of the prescribed rank on any finite-dimensional CW complex? The report says yes through a space-independent universal multiplier, including complexes with infinitely many cells, and gives a telescope example. The cited lemma and its infinite-cell justification require primary review.
\subsection{Sources}
\begin{itemize}
\item Local UnsolvedMath registry, permanent ID $20000250$, source code \texttt{catalogue entry}; created and updated 2026-08-14. Status \texttt{solved}; classification \texttt{SOLVED-IN-LITERATURE}; difficulty L3; proposer field null, while original attribution names Madhav Nori.
\item AIM Problem Lists, Algebraic geometry, workshop \emph{Projective modules and $\A^1$-homotopy theory}, Question 13, source type PDF, \url{https://aimath.org/pastworkshops/projectiveA1problems.pdf}. Canonical location \texttt{aim-algebraic-geometry-notes.json}, \texttt{notes[249]}, tag \texttt{question}; statement status \texttt{exact}. Verification reports direct PDF checking because of severe OCR damage; its coefficient quotient is interpreted as $\Z/m\Z$. Original remarks and literature fields empty; this draft makes no new source-visit claim.
\item The research summary names Benoist's 2026 universal-multiple lemma and target-skeleton proof, without exact title, identifier, theorem statement, or bibliography.
\item Supplied research attempt 1: result \texttt{full\_solution}, run status \texttt{solved}, rigor label \texttt{proved}, low novelty confidence. Candidate contribution \texttt{explicit\_example\_and\_corollary}: telescope bundles and the uniform finite-dimensional corollary. Explicit caveat: machine-generated research aid, not peer reviewed; independently verify before citation.
\item Dataset attribution: Ulam AI / UnsolvedMath Contributors, CC BY 4.0; original problem attribution remains to the supplied AIM workshop collection.
\end{itemize}
\end{document}
